\section{Definicje}
\begin{frame}
  \frametitle{Gra komputerowa}
  \large Aplikacja dostępna na platformie "Steam" oznaczona typem "Game"
\end{frame}
\begin{frame}
\frametitle{Silnik do gier}
\large Oprogramowanie zaprojektowane i stworzone do kreacji gier komputerowych
\end{frame}

\begin{frame}
  \frametitle{Nowoczesne}
  \large Ponad 1000 gier w tej dekadzie na platformie "Steam"
\end{frame}


{
  \setbeamercolor{footline}{fg=white}
\usebackgroundtemplate{\includegraphics[width=\paperwidth, height=\paperheight]{steamdb_main.png}}
\begin{frame}
\end{frame}
}


{
\usebackgroundtemplate{\includegraphics[width=\paperwidth, height=\paperheight]{steamdb_filter.drawio.pdf}}
\begin{frame}
\end{frame}
}

\begin{frame}
  \frametitle{Wybrane silniki - start}
  \begin{center}    
  
\includegraphics[width=0.8\paperwidth, height=0.8\paperheight]{usedEngines.pdf}
\end{center}
\end{frame}

\begin{frame}
  \frametitle{Wybrane silniki}
  \begin{itemize}
    \item Wyeliminowanie nie generycznych - Ren'Py, RPGMaker
    \item Wybór najpopularniejszych - Unity, Unreal 
    % Prawie 25k i ponad 7.5k
  \end{itemize}
\end{frame}

\begin{frame} 
  \frametitle{Wydajność silnika}
  \begin{itemize}
    \item Klatki na sekundę (FPS)
    \item Zużycie CPU, GPU, RAM i VRAM
    \item Liczba draw calls 
    \item Czas ładowania assetów 
    \item Czas odpowiedzi na interakcję gracza
  \end{itemize} 
\end{frame}
\begin{frame}
    \frametitle{Możliwości Silnika}
    \begin{itemize}
        \item Renderowanie grafiki 
        % Ray tracing, HDR lighting, dynamic shadows, particle systems, animacja
        \item Silnik Fizyczny 
        \item Multiplatformowość (VR)
        % Linux, Windows, MacOS, Android, IOs, Xbox, PlayStation, Nintendo, VR
        \item Skryptowanie logiki gier (AI)
        \item Gry online
        \item Sklepy z assetami
    \end{itemize}
  \end{frame}
